We construct the family and verify the class, observed-arm, separation, and information properties.  Begin with \(S=[-1/2,1/2]^2\), lower and upper half-square assignment regions, and
\[
 \mathcal B=\{(u-1/2,0):0\leq u\leq1\},
 \qquad \gamma(u)=(u-1/2,0).
\]
The base density is one.  The chart is injective and smooth of every order.  At an endpoint either observed arm contains a quarter disk.  All densities below are at least \(3/4\), so, for \(0<r\leq r_0\leq1/4\),
\[
 P\{T=t,\ X\in B_2(x,r)\}\geq\frac{3\pi}{16}r^2\geq c_-r^2.
\]
Thus actual-arm thickness holds uniformly.

Choose \(\phi\in C_c^\infty(B_2(0,1))\), \(0\leq\phi\leq1\), with \(\phi(0)=1\), and put
\[
 \Delta=d a_n,
 \qquad w=(d_1\Delta)^{1/q}.
\]
Pack the trimmed interface with centers \(x_1,\ldots,x_M\) separated by at least \(3w\).  For all large \(n\), the radius-\(w\) disks are disjoint and avoid the endpoints, while interval packing gives
\[
 M\geq c w^{-1}\geq c_0a_n^{-1/q}.
\]
Write \(\phi_j(z)=\phi((z-x_j)/w)\).  For \(0\leq k\leq q\),
\[
 \lVert D^k(\Delta\phi_j)\rVert_\infty
 \leq C_\phi\Delta w^{-k}\leq C_\phi d_1^{-1}.
\]
Disjoint supports give the same bound for every signed sum.  Choose \(d_1\) large and an affine slope \(b>0\) small enough for the fixed H\"older ball.

For \(\omega\in\{0,1\}^M\), set
\[
 p_{0,\omega}(z)=\frac12,
 \qquad
 p_{1,\omega}(z)=\frac12+bz_1
 +\Delta\sum_{j=1}^M\omega_j\phi_j(z).
\]
After decreasing \(b\) and \(d\), these lie in \([1/4,3/4]\).  Conditional on \(X=z\), take independent
\[
 Y(t)=B_t+\varepsilon_t,
 \qquad B_t\sim\operatorname{Bernoulli}(p_{t,\omega}(z)),
 \qquad \varepsilon_t\sim N(0,1).
\]
Then \(\mu_{t,P_\omega}=p_{t,\omega}\) and \(1\leq\sigma^2_{t,P_\omega}\leq5/4\).  Hoeffding's lemma and the Gaussian moment-generating identity yield
\[
 \mathbb E\!\left[e^{\lambda\{Y(t)-p_{t,\omega}(X)\}}\mid X\right]
 \leq e^{\lambda^2/8}e^{\lambda^2/2}
 \leq e^{L^2\lambda^2/2}.
\]

To hide the bump at its center, choose a continuous cutoff \(\chi\) with \(\chi(v)=0\) for \(v\leq1\) and \(\chi(v)=1\) for \(v\geq2\), and set
\[
 A_n(r)=-\frac{2\Delta\phi(r/w)}{br}
 \chi\!\left(\frac{br}{8\Delta}\right).
\]
There is no division at zero because the cutoff vanishes near zero.  On the active region \(br\geq8\Delta\), \(|A_n(r)|\leq1/4\) after enlarging the fixed cutoff constant.  Smooth compact support makes the correction continuous at its outer edge.  Put
\[
 f_\omega(z)=1+\sum_{j=1}^M\omega_j
 A_n(\lVert z-x_j\rVert_2)
 \frac{z_1-(x_j)_1}{\lVert z-x_j\rVert_2}
 \mathbf1\{z\in B_2(x_j,w)\},
\]
with every summand defined as zero at its center.  Each angular correction integrates to zero on every full circle, so \(\int_Sf_\omega=1\).  Disjointness makes the formula unambiguous, continuous, and bounded in \([3/4,5/4]\subseteq[L^{-1},L]\).  Set \(T=\mathbf1\{X\in\mathcal E_{1,P_\omega}\}\), \(Y=TY(1)+(1-T)Y(0)\), and sample independently.  This verifies assignment, consistency, iid sampling, compact support, absolute continuity, the density envelopes, H\"older smoothness, sub-Gaussian tails, and actual-arm thickness for every vertex.

Compare adjacent vertices \(\omega\) and \(\omega^{(j)}\), oriented so that \(\omega_j=1\).  On either observed half-circle centered at \(x_j\), the angular integral of \(\cos\theta\) is zero.  Thus their observed arm-specific radial measures coincide, and their fixed RN densities coincide Lebesgue-almost everywhere.  On the treated upper half-circle, whenever \(br\geq16\Delta\),
\[
\begin{aligned}
 &\int_0^\pi p_{1,\omega}(x_j+re_\theta)
 \{1+A_n(r)\cos\theta\}\,d\theta
 -\int_0^\pi p_{1,\omega^{(j)}}(x_j+re_\theta)\,d\theta\\
 &\qquad=\pi\Delta\phi(r/w)+\frac\pi2brA_n(r)=0.
\end{aligned}
\]
The control equality follows from its constant regression.  A Bernoulli variable followed by the common Gaussian channel is determined, after angular mixing at a fixed radius, by the averaged Bernoulli probability.  Hence the observed radial outcome laws agree outside \(r\leq16\Delta/b\), modulo the immaterial null set of RN representatives.

Inside that interval the common radial density is at most \(Cr\), and the averaged Bernoulli probabilities differ by at most \(C\Delta\) while remaining in \([1/4,3/4]\).  The inequalities
\[
 \operatorname{kl}(p,p')\leq\frac{(p-p')^2}{p'(1-p')},
 \qquad
 \operatorname{KL}(B+\varepsilon,B'+\varepsilon)
 \leq\operatorname{kl}(p,p')
\]
and data processing therefore give
\[
 \operatorname{KL}(P^{\mathrm{dist}}_{\omega,x_j},
 P^{\mathrm{dist}}_{\omega^{(j)},x_j})
 \leq C\int_0^{16\Delta/b}\Delta^2r\,dr
 \leq C'\Delta^4.
\]
Tensorization yields \(C'n\Delta^4=C'd^4\log n\), whereas
\[
 \log M=\left(\frac1{4q}+o(1)\right)\log n.
\]
Choose \(d\) small enough that the divergence is at most \(\alpha\log M\) for a fixed \(\alpha<1/8\), retaining slack for Poissonization at mean \(2n\).  At \(x_j\), adjacent targets differ by \(\Delta\phi_j(x_j)=da_n\), so \(c_1=d\) works.

For the canonical square \([-1,1]^2\), use \(\gamma(u)=(2u-1,0)\), base density \(1/4\), and the same relative-density construction.  If \(L>4\), put \(\varepsilon_L=1-4/L>0\), enlarge the cutoff constant depending only on \(L\), and choose \(d_1\) large and then \(d/b\) small.  The relative perturbation is at most \(\varepsilon_L/2\), so
\[
 L^{-1}\leq\frac14(1-\varepsilon_L/2)
 \leq f_\omega\leq\frac14(1+\varepsilon_L/2)\leq L.
\]
All cancellation and divergence bounds change only by constants, placing the family in \(\mathcal P_{\mathrm{straight}}\).  At \(L=4\), normalization and the floor force density \(1/4\), so the nonzero angular perturbation is unavailable; no such construction is claimed there.